\subsection{Agreement with measured target relationships improves on some panels}

We compared the raw and residual maps with target relationships measured in external assays.
The kinase endpoint comes from the experimental panels alone and uses no docking score. On the
20 kinases shared by DAVIS,
PKIS2 and PKIS1, a target pair counts as positive when its centred experimental correlation
falls in the upper decile of at least two of those panels, which labels 15 of the 190 pairs;
we refer to this as the frozen analysis endpoint. The release does not establish when it was fixed
relative to inspection of KiRHub. Inference permutes the 20 docking target labels 50{,}000 times
while the experimental endpoint stays fixed.

Raw Vina target geometry gives an area under the receiver operating characteristic curve
(ROC-AUC) of 0.597 and average precision of 0.236. Removing each ligand's mean raises these to
0.710 and 0.301, respectively (Figure~\ref{fig:external}b). The paired increments are $+0.114$
and $+0.066$. These gains are descriptive: the analysis is post hoc, and the result does not
remain below the 0.05 threshold when the full set of cutoffs is considered.

Two structural baselines on the identical 190 pairs set the scale. Receptor-domain sequence
identity gives ROC-AUC 0.718 and average precision 0.426, and identity across the aligned
85-residue KLIFS ATP pocket gives 0.708 and 0.486. On ROC-AUC, the centred map approaches these
structural baselines: it is 0.008 below sequence identity and
0.002 above pocket identity, having started about 0.12 below each. These point estimates do not
establish equivalence or superiority. On average precision it remains below
both. The corresponding gains in continuous Spearman agreement over the same 190 pairs are much
smaller: $+0.016$ on DAVIS and $+0.044$ on PKIS2.

PDSP gives the largest point-estimate separation between the two docking maps, but it also shows
the clearest support boundary. Its exact-only graph contains 148 target pairs over 21 targets,
estimated from 10 to 92 shared compounds per edge. Before building the docking map, we removed every
Docking-44 row whose connectivity block appeared among any analyzed PDSP compound. The exact-only
endpoint shared 63 full keys and 116 connectivity blocks with Docking-44; the broader conservative
mask removed 128 docking rows. Raw-map
agreement is $\rho=0.044$, whereas residual-map agreement is 0.295. The smaller Certified Data
subset shows the same direction, 0.074 versus 0.425 (Figure~\ref{fig:external}a).

Two-sided paired QAP places this separation in the tail under unrestricted target relabelling:
the raw-to-residual difference in continuous agreement is 0.251, with $p=0.0020$ and Holm-adjusted
$p=0.0240$ across the 12-row table. Under within-family relabelling, the same descriptive contrast
gives $p=0.0134$ and Holm-adjusted $p=0.1016$. Adding residual docking to a sequence-and-family
baseline raises continuous agreement by 0.096 on this support, but none of the adjusted fusion tests
falls below 0.05. These probabilities describe a post-hoc analysis family, not a confirmatory test.

The support sensitivities are more important than the smallest probability. Restoring right-censored
cells at their reported bounds lowers residual agreement to 0.076. Rebuilding both docking maps on
the 9{,}221 rows complete across all 44 targets changes the comparison to $-0.099$ for raw docking
and 0.046 for residual docking; the sequence-and-family fusion then becomes slightly worse when
residual docking is added. This deletion is driven almost entirely by missing scores for CHRM2.
Complete-case deletion changes the chemical domain, so this is a
support-selection sensitivity rather than a pure imputation contrast. Raising the minimum shared-compound
rule from 10 to 30 preserves the direction but changes the graph size and effect magnitude. QAP keeps
edges that share a target together, but does not propagate uncertainty from estimating each edge on
its own compound set. We therefore treat the PDSP result as descriptive and support-sensitive.

Not every receptor in the panel is docked into a human structure, and the sequence baseline is
human throughout. An audit of RCSB PDB metadata for all 43 accessions \cite{berman2000} finds five non-human receptor chains:
GRIN1 (rat), ADRB1 (turkey, eight mutations), PTGS1 (sheep), PTGS2 and HTR3A (mouse); the other
non-human chains in these entries are crystallisation aids on a human receptor --- T4 lysozyme
and BRIL fusions, nanobodies, murine Fab fragments, G-protein and scFv16 chains --- and are not
species substitutions
(\path{results/docking_structure_species/structure_species.csv}). Two of the five, ADRB1 and
HTR3A, fall in the PDSP graph. Dropping them leaves 119 pairs over 19 targets, on which the
increment over the sequence-and-family baseline is $+0.061$ in continuous agreement
and $+0.063$ in top-decile retrieval. The point estimates are smaller on this construct-clean sensitivity but
remain positive after removing the two proxied receptors.

The Novartis safety panel provides an exact- and connectivity-de-overlapped test. It covers 12 targets and
1{,}975 InChIKeys, and all shared connectivity blocks were removed before the docking map was
built. Most selected activity rows are right-censored. Against a fixed centred assay map, the
reported-bound endpoint moves from $\rho=0.149$ to 0.372; the 10\,$\mu$M endpoint changes only
from 0.348 to 0.388. On the largest outcome-blind complete rectangle, 102 compounds measured on
the same 11 targets, agreement moves from 0.055 to 0.575 (Figure~\ref{fig:external}a), and all
leave-one-target-out differences on the reported-bound endpoint are positive. The remaining
endpoint and assay-transform values are in Supplementary Section S6.

After rank control for sequence identity, family, pair support and target coverage, only one of
three residual associations has a one-sided family-preserving tail area below 0.05. The raw map
also retains partial agreement and exceeds the residual map in partial rank at 10\,$\mu$M, so this does not
establish a general increment beyond homology. On the complete 102-compound rectangle, residual
partial rank is $+0.371$, compared with $-0.220$ for the raw map.

Raw and centred experimental maps encode different edges: the former captures shared response
magnitude, whereas the latter captures target-relative response. Changing both the assay and docking
transforms at once would therefore make the docking-side effect uninterpretable. Every comparison in
Figure~\ref{fig:external} holds the experimental map fixed and varies only the docking map.

On KiRHub the comparison moves in the opposite direction. Against the upper decile of the
centred KiRHub target-pair map, 19 positives among the same 190 pairs, raw docking gives
ROC-AUC 0.751 and average precision 0.331, while centred docking gives 0.744 and 0.289. The
paired target-label probabilities for a positive gain are 0.53 on ROC-AUC and 0.92 on average
precision. This is a non-replication of the centred-minus-raw gain on a fourth kinase panel,
showing that the sign of the docking-side effect is not fixed across panels. The raw and centred
maps are nevertheless each associated with the experimental map on their own
(ROC-AUC $p=0.0063$ raw and $p=0.0024$ centred).
